\chapter{The phases: what each proves, and how}\label{ch:phases}

This chapter gives, phase by phase, the design as the blueprint has it, the design as leanVM
has it, and the intuition behind the two proofs each phase owes: perfect completeness (an
honest prover always passes) and round-by-round knowledge soundness (an invariant on partial
transcripts that a cheating prover escapes only when a fresh challenge lands badly, with a
bound per challenge). Where the blueprint falls short of leanVM the finding is named and
detailed in \cref{ch:faithfulness}; where a definition cannot be met or says less than it
should, in \cref{ch:nonvacuity}. The vocabulary of the ideal model (oracle reduction,
knowledge state function, extractor) is introduced in \cref{ch:libraries}; the reader who has
that chapter in mind can read this one on its own.

\section{The commit phase}\label{sec:ph-commit}

The prover sends the stack $q$ as the one oracle message; the verifier keeps the statement and
exposes the message as the oracle. There is no challenge and no check. The phase is the spine's
own and both its halves are proved (\lean{commitComplete}, \lean{commitSecurity}). Its knowledge
state function says, before the message, that the candidate witness satisfies \lean{M3Holds}
and, after it, that the \emph{message} does; its extractor reads the message. Error zero is
right: there is no challenge to charge. In the deployed protocol the message is a Merkle root;
what a root must provide for the compiled protocol to inherit the oracle protocol's guarantee
is the subject of \cref{sec:ph-opening}.

\section{The bus phase}\label{sec:ph-bus}

\paragraph{What is proved.} That the multiset of tuples the tables push equals the multiset
they pull, on each of the three buses, and that every read count is nonzero. With the boundary
blocks (the initial and final machine state, the memory and bytecode seeds and finalizations)
counted in, balance is the whole of the memory and bytecode consistency of the execution.

\paragraph{How.} A fingerprint turns a 16-tuple $t$ into the scalar
$\beta + \sum_{i < 16} \eq(\alpha, i)\, t_i$ for random $\alpha \in \E^4$, $\beta \in \E$; two
multisets of at most $2^{\mu_{\mathrm{bus}}}$ tuples with different contents have different
products of fingerprints except with probability $4 \cdot 2^{\mu_{\mathrm{bus}}} / \abs{\E}$
(the specification's Theorem 5.1; each factor has total degree 4 in $(\alpha, \beta)$). The
product itself is proved by a grand-product argument: the leaves are the fingerprints (padded
with $1$), each layer of the tree holds pairwise products, and a GKR walk from the root
descends layer by layer, reducing a claim on a layer to a claim on the layer below by a
sumcheck; the three trees (push, pull, count) share every challenge and are batched by the
powers of a combiner $\lambda$. The prover states a single root $R$ for the push and pull
trees, so their equality needs no check, and a root $R_c$ for the count tree, which the
verifier requires to be nonzero. The walk ends with three claims on the leaf tables at one point
$\zeta$, which the leaf decomposition (the specification's equation (2) of §5.4) splits into a
public part the verifier computes, five boundary evaluations the prover sends, and the tables'
share, which the table sumcheck settles.

\paragraph{Completeness.} From \lean{M3Holds}: the constraints vanish, so the zerocheck claims
hold at every point; balance is a permutation, so the two products are one field element and
the honest prover has a single $R$ to send; every count is nonzero, so $R_c \neq 0$; the GKR's
honest messages pass every layer check (no inverse of a challenge is taken); the leaf
decomposition holds of the honest stack with the derived $\mathit{rem}_s$. This is perfect.

\paragraph{Knowledge soundness.} One state function for the whole phase, a conjunction:
before $(\alpha, \beta)$, \lean{M3Holds} of the oracle; after them, \enquote{the two products at
$(\alpha, \beta)$ are equal and $R_c$ is the true count product and every constraint vanishes};
from then on the GKR's own invariant (\enquote{the current layer claims are the true values})
together with the zerocheck conjunct \enquote{every constraint's extension, restricted to the
coordinates of $\zeta$ drawn so far, is identically zero}. A challenge flips a false state
to true with probability at most: $4 \cdot 2^{\mu_{\mathrm{bus}}} / \abs{\E}$ for
$(\alpha, \beta)$; $2 / \abs{\E}$ for a combiner (the batching polynomial has degree 2);
$4 / \abs{\E}$ for a round challenge (the cofactor has degree 4); $1 / \abs{\E}$ for a
combination challenge; and $0$ for the last, unused combiner. The zerocheck conjunct adds
nothing: on a coordinate of $\zeta$ it flips with probability at most $1 / \abs{\E}$, once for
all constraints (a nonzero multilinear restricted at one more variable becomes zero for at most
one value of that variable), and a conjunction escapes with the larger of the two
probabilities, not their sum. The blueprint's three different accounts of this charge, and its
GKR of the wrong degree, are the findings of \cref{sec:faith-bus}.

\paragraph{Where the blueprint falls short.} The GKR layer sumcheck it specifies has degree 5
and is not the deployed one (\emph{normalized} cofactor of degree 4, no equality factor in the
layer check); it draws one combiner fewer than the deployed verifier; its generic GKR takes the
leaf tables as oracles, which the spine's one-oracle phases cannot supply without ArkLib's
admitted context lifting; and the zerocheck conjunct, which changes on the GKR's own
challenges, cannot be framed around a GKR that knows nothing of the constraints. The admissible
sizes it equates with leanISA's caps leave out the stacking window and the rate window that the
deployed verifier checks. None of this touches the spine as built: the real bus phase's output
is expressible as a \lean{BusOut}, and its completeness and knowledge soundness against the
spine's two seams can hold (\cref{sec:faith-bus}).

\section{The table sumcheck}\label{sec:ph-table}

\paragraph{What is proved.} That every constraint of every opcode table vanishes on every row
(for leanVM at the pin, the two identities of the JUMP table; the other tables have none), and
that each table's bus forms sum, over its rows, to the share $\mathit{rem}_s$ the bus phase
derived for it. The second is what ties the committed columns to the bus trees: the trees were
built by the prover, and only this sumcheck says that they were built from the columns.

\paragraph{How.} One sumcheck over $\tau_{\max}$ variables, the highest bound first, with each
table joining at the round where its height is reached and padded below it by the product of
the variables it lacks (so that its sum is unchanged). The summand batches the constraints and
the three bus forms by the powers of $\xi$, with $\eq(\zeta_{< \tau_j}, \cdot)$ as the weight
of table $j$: the constraint terms contribute zero to the target and the bus terms contribute
$\sum_s \xi^{2 + s}\, \mathit{rem}_s$, which the verifier computes, so no target is ever sent.
A round message is a cubic (the equality factor times a degree-2 term) of which three
coefficients travel; the verifier derives the fourth from the running claim. At the end the
prover sends every column's value at the sumcheck's point and the verifier recomputes the
summand.

\paragraph{Completeness.} Immediate from the input seam: the zerocheck claims are true, the bus
forms sum to their values, the honest round polynomials pass, the honest column values
reproduce the summand. No inverse of a challenge is taken.

\paragraph{Knowledge soundness.} Before $\xi$: every claim is true. After $\xi$: the batched
claim $\sum_c \xi^c\, \delta_c = 0$ where $\delta_c$ is claim $c$'s error; a nonzero polynomial
of degree at most $k - 1$ in $\xi$ has at most $k - 1$ roots, so the error on $\xi$ is
$(k - 1) / \abs{\E}$, which is $(B + 2) / \abs{\E}$ for $B$ constraints and three sides. Each
round: if the running claim is false, the prover's cubic differs from the true one and they
agree at most at three points, $3 / \abs{\E}$. After the last message: the verifier accepts and
every column claim holds, so the recomputed summand is the true summand at the point and the
running claim was true. The escape from \enquote{the zerocheck claims hold at $\zeta$} to
\enquote{the constraints vanish on the cube} belongs to the bus phase, where $\zeta$ was drawn
(above); it is not charged here.

\paragraph{Where the blueprint falls short.} Not in the numbers, which are right and tight,
but in the seam it consumes. \lean{Seam.bus} as built admits statements that the deployed table
sumcheck cannot serve: any number of linear claims (so the error on $\xi$ is unbounded over the
seam), terms of one table at unrelated points, terms on any table of the instance, and a degree
clause on the statement alone, which forces the verifier to check degrees, a check leanVM does
not make. Completeness and knowledge soundness are owed on all of them. The blueprint also
leaves open which tables the sumcheck ranges over: its leanISA instance makes the six shared
columns tables of the instance, and read literally its table sumcheck then has at least sixteen
rounds where leanVM has as few as three (\cref{sec:faith-table}). And its oracle protocol sends
four coefficients and checks the round identity, where the deployed verifier reads three and
derives the fourth; the lemma relating the two is not stated.

\section{The public-input phase}\label{sec:ph-pub}

\paragraph{What is proved.} That cells 0 and 1 of the three memory limbs hold the two public
words: for the two low limbs, the words' limbs; for the top limb, zero.

\paragraph{How.} The verifier draws $r \in \E$ and considers the line through cells 0 and 1 of
a column: $\widetilde{c}(r, 0, \dots, 0) = (1 - r)\, c(0) + r\, c(1)$, where $-1 = 1$ in
characteristic 2. The prover sends the values $c_0, c_1$ it claims for the two low limbs on
that line. The specification checks each against the line through the public words' limbs;
the deployed verifiers check the single equation $c_0 + y\, c_1 = (1 + r)\, w_0 + r\, w_1$ in
$\E$, which the two per-limb equations imply and which does not imply them. Three claims are
then pooled: the two low limbs at the values sent and the top limb at $0$, for which nothing is
sent.

\paragraph{Completeness.} If the lines hold, the honest values pass either check and every
pooled claim is a true evaluation. Perfect.

\paragraph{Knowledge soundness.} A wrong stack (a cell of a limb differing from the public
word's) makes the pooled claim on that limb a nonzero polynomial of degree one in $r$, true at
one challenge at most: error $1 / \abs{\E}$, whatever the number of lines and for either
check. For the deployed check the argument runs through the independence of $1, y, y^2$ over
$\K$: if the three pooled claims hold of a $\K$-valued memory and the combined equation holds
at two distinct challenges, the cells are the limbs of the public words and the top limbs are
zero. This is proved in Lean by the review (\cref{ch:probes}).

\paragraph{Where the phase as built falls short.} It follows the specification, so its theorem
is about a verifier that none of the three executable verifiers runs. And it pools the values
it \emph{computes} from the statement rather than the values \emph{sent}: on the accepting path
they agree, but the consequence is that the check on the prover's message is not load-bearing
for any theorem of the phase: a verifier that reads the message and ignores it has the same
completeness and the same knowledge soundness at the same error. The requirement that a missing
check leave a theorem unprovable is therefore not met by this phase as designed, and cannot be;
pooling the values sent, as the deployed verifiers do, restores it (\cref{sec:faith-pub},
\cref{ch:drift}). The top-limb claim exists because the leanISA instance lists a third public
line; nothing in the phase or the spine forces that line, and what catches its absence is the
adaptor's theorem, not built (\cref{ch:drift}).

\section{BLAKE2s validity and ring switching}\label{sec:ph-flock}

\paragraph{What is proved.} That the bits packed into the region $\qflock$ of the stack satisfy,
block by block, the R1CS of the BLAKE2s compression with the constant position of every block
equal to $1$; the eighteen value limbs each BLAKE2S table row carries are strided slots of that
region and are opened with every other point claim.

\paragraph{How.} Flock is a zerocheck over the R1CS's rows, with the first six variables
skipped univariately (the prover sends the summand on a coset of 64 points), followed by a
lincheck that reduces the three matrix products to one evaluation of the bit witness at a
point, with a single terminal identity the verifier checks. Ring switching then turns the 64
resulting claims on the bit polynomials $Q_i$ (bit $i$ of every word) at one point into a
weighted claim on the packed column: for a random $\F_2$-linear map $\Phi$ built from six
challenges, $\sum_u \Phi(\eq(r, u))\, \qflock(u) = \sum_i x^i\, \Phi(s_i)$, whose weight the
verifier evaluates in closed form. Neither step commits anything: the only oracle is $q$.

\paragraph{Completeness.} Perfect, for the honest prover that sends the coefficients of the
true round polynomials: the verifier derives the untransmitted coefficient by a multiplication
and takes no inverse. The \emph{deployed} Rust prover is not that prover: it derives $G(0)$
through $(1 + r_{\mathrm{eq}})^{-1}$, whose value at $r_{\mathrm{eq}} = 1$ is $0$ by the
field's convention, and at that challenge emits a proof its own verifier rejects. The
blueprint's acceptance test 20 attributes this inverse to the protocol and files it under a
soundness error; both are wrong (\cref{sec:faith-flock}).

\paragraph{Knowledge soundness.} The specification's Annex C accounts
$(4 k_{\mathrm{batch}} + 163) / \abs{\E}$ for the reduction, which the review re-derived term by
term from the Rust's parameters, and Annex A accounts below $2^{32} / \abs{\E}$ for ring
switching (a nonzero polynomial of total degree below $2^{32}$ in the six challenges). Both are
errors for a fixed committed polynomial; after compilation the list-decoding regime multiplies
them by the list size (\cref{sec:ph-opening}).

\paragraph{Where the blueprint falls short.} Its Layer 9 makes the eighteen limb claims an
\emph{input} of Flock and one weighted claim its only output, dropping the limb claims; with
leanVM's verifier, which reads no claim, that interface has knowledge error one. The spine's
slot, which keeps the column claims and adds the weighted one, is right. The slot map of the
limbs is placed inside the Flock interface while the instance's layout needs it. And a Flock
phase cannot be written over an abstract instance at all, since the instance offers it only an
opaque predicate \lean{aux}; the review proposes a small structure that names the region, the
number of compressions and what \lean{aux} says of them (\cref{sec:faith-flock}). Today, with
the only built instance having \lean{aux} trivially true, a Flock phase that checks nothing is
knowledge sound at error zero (\cref{ch:probes}): the master theorems constrain the Flock
verifier exactly as much as the instance's predicate does.

\section{The opening and the compilation}\label{sec:ph-opening}

\emph{This section is completed from the dossier on the opening, WHIR, Merkle trees and
Fiat--Shamir, which is being finished as this draft is written.} What is already established:
the specification's opening is one batching challenge followed by one inner-product opening,
with no separate reduction sumcheck, whereas the blueprint's oracle protocol runs a sumcheck
reducing the batched weighted claim to an evaluation and then queries the oracle once; the
oracle interface of the stack in Layer 0 is an evaluation oracle, and an inner-product oracle,
the idealization of the specification's Definition 3.13, would make the opening phase
\enquote{$\lambda$, then one query} with leanVM's transcript. The compiled error the blueprint
sketches lacks the hash-collision term, the list-size factor of a list-binding commitment and
any account of grinding; the assumed Fiat--Shamir and BCS interfaces are stated for
constructions that are not leanVM's single BLAKE2s map (\cref{sec:faith-opening}).
